\documentclass[11pt,a4paper]{article}
\usepackage[T1]{fontenc}
\usepackage[utf8]{inputenc}
\usepackage{lmodern,amsmath,amssymb}
\usepackage[margin=25mm]{geometry}
\usepackage[hidelinks]{hyperref}
\hypersetup{pdftitle={The map as one conditional theorem: two hypotheses, both finite in kind, and an explicit l},pdfauthor={navstokgap research working notes}}
\newcommand{\tightlist}{\setlength{\itemsep}{0pt}\setlength{\parskip}{0pt}}
\setlength{\emergencystretch}{3em}
\setcounter{secnumdepth}{0}
\begin{document}
\hypertarget{the-map-as-one-conditional-theorem-two-hypotheses-both-finite-in-kind-and-an-explicit-lower-bound-mgehbar-cgammaa_}{%
\section{\texorpdfstring{The map as one conditional theorem: two
hypotheses, both finite in kind, and an explicit lower bound
\(m\ge\hbar c\,\gamma'/a_*\)}{The map as one conditional theorem: two hypotheses, both finite in kind, and an explicit lower bound m\textbackslash ge\textbackslash hbar c\textbackslash,\textbackslash gamma\textquotesingle/a\_*}}\label{the-map-as-one-conditional-theorem-two-hypotheses-both-finite-in-kind-and-an-explicit-lower-bound-mgehbar-cgammaa_}}

Everything this programme has established about the \(SU(3)\) mass gap
can be stated as a single conditional theorem. Its two hypotheses are
the two things not proved here: control of the blocking steps from the
weak side down to the coupling where the correlation length is one
lattice unit, and a certified mixing statement at that one coupling on
one finite box. Its conclusion is the Jaffe--Witten gap along the
trajectory with an explicit lower bound,
\[m\ \ge\ \frac{\hbar c\,\gamma'}{a_*},\] where \(a_*\) is the physical
lattice spacing at which the trajectory reaches the verified coupling,
about \(0.1\) to \(0.17\) fm by the published scale, and \(\gamma'\) is
the certified decay rate per lattice unit there. Any positive
\(\gamma'\) gives a positive gap, and the data say the true rate is
about \(0.8\). Constants explicit; nothing promoted.

\hypertarget{setting}{%
\subsection{1. Setting}\label{setting}}

Wilson measure \(\mu_{\beta_0}\) on the links of \((a_0\mathbb Z)^4\) at
bare coupling \(\beta_0=6/g_0^2\), in a periodic box of any size. A
block-spin map \(\mathcal B\) of block size \(2\): a gauge-covariant
assignment of a coarse link variable to each block, measurable in the
fine links, such that coarse Wilson loops are functions of fine Wilson
loops. The renormalized measures \(\mu^{(k)}=\mathcal B^k\mu_{\beta_0}\)
live on the links of \((2^ka_0\mathbb Z)^4\) and, by construction,
reproduce exactly the correlations of coarse observables. Write
\(\Phi^{(k)}\) for the renormalized interaction when it exists, and
\(\|\Phi\|_\kappa=\sup_\ell\sum_{X\ni\ell}e^{\kappa\operatorname{diam}X}\|\Phi_X\|_\infty\)
for the interaction norm with range weight \(\kappa>0\).

\hypertarget{the-hypotheses}{%
\subsection{2. The hypotheses}\label{the-hypotheses}}

\textbf{H1 (the weak side).} There is \(K=K(a_0)\) with
\(2^Ka_0\to a_*\) as \(a_0\to0\), a rate \(\gamma_*>0\) and constants
\(C_k\) such that for each \(k<K\) and all local observables \(F,G\) of
the \(k\)-th lattice with supports at distance \(d\) in its units,
\[\Big|\langle F;G\rangle_{\mu^{(k)}}-\langle\bar F;\bar G\rangle_{\mu^{(k+1)}}\Big|
\ \le\ C_k\,\|F\|_\infty\|G\|_\infty\,e^{-\gamma_*d},\] where \(\bar F\)
is the conditional expectation of \(F\) given the coarse links; and
\(\Phi^{(K)}\) exists, is finite in \(\|\cdot\|_\kappa\), and satisfies
\(\|\Phi^{(K)}-\Phi_0\|_\kappa\le r\) for a reference interaction
\(\Phi_0\) and radius \(r\) fixed in H2.

In words: at every step before the last, the fluctuations integrated out
cluster exponentially at the scale of the step, and the trajectory
arrives within a fixed distance of a reference interaction. For the
steps at which the effective coupling is small this is the content of
the constructive programme, perturbative in kind and unavailable with
constants; for the last three, in which \(\xi/a\) passes from about
\(10\) to about \(1\), it is the open problem
(\href{confinement-scale-bands.md}{bands note} §2b).

\textbf{H2 (one box).} The reference interaction \(\Phi_0\) satisfies
the strong-mixing condition on a box \(V\) of side \(R\): for every
\(W\subseteq V\), every boundary link \(y\) and boundary conditions
\(\omega,\omega'\) differing at \(y\),
\[\big\|\mu_V^{\omega}\big|_W-\mu_V^{\omega'}\big|_W\big\|_{\rm TV}\ \le\ C\,|W|\,e^{-\gamma\,d(W,y)},\]
with a margin such that every \(\Phi\) with
\(\|\Phi-\Phi_0\|_\kappa\le r\) satisfies the same condition with some
\(C',\gamma'>0\). The natural \(\Phi_0\) is the Wilson interaction at
\(\beta_W\simeq6\), where the published correlation length is about one
lattice unit, or, better, a nearby interaction with a small negative
adjoint term, which moves \(\Phi_0\) away from the endpoint of the
fundamental--adjoint first-order line where the scalar channel softens
(precision of 2026-09-23, \href{mass-gap-openings.md}{openings note}
§1); \(R\) is \(3\) to \(5\); the condition is to be read on the
gauge-invariant content of \(W\), the small Wilson loops, since
single-link marginals of interior links are Haar
(\href{confinement-scale-bands.md}{bands note}).

\hypertarget{the-theorem}{%
\subsection{3. The theorem}\label{the-theorem}}

\textbf{Theorem (conditional).} Assume H1 and H2. Then for every bare
coupling on the trajectory and every periodic volume:

\begin{enumerate}
\def\labelenumi{\arabic{enumi}.}
\tightlist
\item
  \(\mu^{(K)}\) has a unique Gibbs state and truncated correlations of
  all local observables decaying at rate at least \(\gamma'\) per unit
  of \(2^Ka_0\) (Dobrushin--Shlosman; Martinelli--Olivieri);
\item
  \(\mu_{\beta_0}\) has truncated correlations of all local fine
  observables decaying at rate at least
  \(\min\big(\gamma_*,\,\gamma'\big)/2^K\) per unit of \(a_0\);
\item
  the Wilson transfer matrix at spacing \(a_0\) has a spectral gap
  \[\Delta(a_0)\ \ge\ \frac{\hbar c}{a_0}\cdot\frac{\min(\gamma_*,\gamma')}{2^{K(a_0)}}\]
  on the cyclic subspace of local observables, uniformly in the volume;
\item
  in the continuum limit along the trajectory,
  \[m\ \ge\ \frac{\hbar c\,\min(\gamma_*,\gamma')}{a_*}\ >\ 0 .\]
\end{enumerate}

\emph{Proof.} (1) is the mixing-to-decay theorem applied to
\(\Phi^{(K)}\), which lies in the verified neighbourhood by H1 and H2.
For (2), write a fine observable as its coarse part plus fluctuation
parts step by step, \(F=\bar F^{(K)}+\sum_{k<K}(F^{(k)}-\bar F^{(k)})\);
the coarse-part correlations decay by (1) at rate \(\gamma'\) per coarse
unit, that is \(\gamma'/2^K\) per fine unit, and the fluctuation-part
correlations at step \(k\) decay by H1 at rate \(\gamma_*\) per unit of
\(2^ka_0\), that is \(\gamma_*/2^k\ge\gamma_*/2^K\) per fine unit; the
cross terms are bounded by the same rates. (3) is link reflection
positivity together with the spectral argument of
\href{wilson-strong-coupling-explicit.md}{the Wilson note} §3: decay at
rate \(\mu\) per fine unit for all local observables forces the spectral
measure of every local vector orthogonal to the vacuum to vanish above
\(e^{-\mu}\). (4) follows from (3) as \(a_0\to0\) with
\(2^{K(a_0)}a_0\to a_*\). \(\square\)

The Jaffe--Witten clause \(0<m<\infty\) is then (4) together with the
upper side of \href{moment-hierarchy-upper-bounds.md}{the
moment-hierarchy note}; T3, the ratio \(m/(\hbar c\Lambda)\), is the
statement that \(a_*\Lambda\) has a limit, which is the convergence of
the trajectory to one curve as in
\href{intermediate-region-finite-verification.md}{the
finite-verification note}.

\hypertarget{what-the-theorem-makes-visible}{%
\subsection{4. What the theorem makes
visible}\label{what-the-theorem-makes-visible}}

\emph{The bound is a ratio of a certified rate to a physical length.}
With \(a_*\simeq0.093\) fm at \(\beta_W=6\) and the measured glueball
mass \(1.66\,\hbar c\,{\rm fm}^{-1}\), the true rate is
\(m a_*\simeq0.78\); a certified \(\gamma'\) of any size gives a
theorem, and one of order \(0.1\) would already be a physically
meaningful bound.

\emph{Both hypotheses are finite in kind.} H2 is one box at one
coupling. H1 is a sequence of steps, uniform in the perturbative regime,
whose non-perturbative part is three steps. Neither is available today:
H2 needs a certified computation, and H1's last three steps are the
confinement-scale problem.

\emph{Nothing else is missing.} T1, T2, the strong-coupling side below
the verified coupling, the transfer from decay to gap, the upper side of
the gap, and the reduction of T3 are all in place with explicit
constants. The theorem is the repository's formal endpoint: a proof of
the mass gap for \(SU(3)\) consists of H1 and H2.

\hypertarget{consequence-for-state}{%
\subsection{5. Consequence for STATE}\label{consequence-for-state}}

The map is a theorem with two hypotheses and an explicit conclusion
\(m\ge\hbar c\min(\gamma_*,\gamma')/a_*\). Progress from here is
progress on H1 or H2, and nothing else.
\end{document}
